% ભાગ: first-order-logic
% પ્રકરણ: models-theories
% વિભાગ: introduction

\documentclass[../../../include/open-logic-section]{subfiles}

\begin{document}

\olfileid{fol}{mat}{int}
\olsection{પરિચય}

\begin{explain}
સ્વયંસિદ્ધિમૂલક પદ્ધતિનો વિકાસ વિજ્ઞાનના ઇતિહાસની નોંધપાત્ર સિદ્ધિ
છે અને ગણિતના ઇતિહાસમાં તેનું વિશેષ મહત્ત્વ છે. કોઈ ક્ષેત્રના
સ્વયંસિદ્ધિમૂલક વિકાસમાં ઘણા પ્રશ્નો સ્પષ્ટ કરવા પડે છે: આ ક્ષેત્ર
શેના વિશે છે? તેના સૌથી પાયાના ખ્યાલો કયા છે? તેઓ પરસ્પર કેવી રીતે
સંબંધિત છે? શું ક્ષેત્રના બધા ખ્યાલોને આ પાયાના ખ્યાલોના પદોમાં
વ્યાખ્યાયિત કરી શકાય? આ ખ્યાલો કયા નિયમોનું પાલન કરે છે અને અનિવાર્યપણે
કરવું જોઈએ?

સ્વયંસિદ્ધિમૂલક પદ્ધતિ અને તર્કશાસ્ત્ર જાણે એકબીજા માટે જ બનેલાં છે.
વિધિવત્ તર્કશાસ્ત્ર સ્વયંસિદ્ધિમૂલક સિદ્ધાંતો ઘડવા, સિદ્ધાંતનાં
સ્વયંસિદ્ધિઓમાંથી ચોક્કસ નિર્ધારિત રીતે પ્રમેયો સાબિત કરવા અને
સ્વયંસિદ્ધિઓને સંતોષતી બધી પ્રણાલીઓના ગુણધર્મોનો વ્યવસ્થિત અભ્યાસ
કરવા માટેનાં સાધનો આપે છે.
\end{explain}

\begin{defn}
!!{sentence}sનો ગણ~$\Gamma$ \emph{સંવૃત} ત્યારે અને માત્ર ત્યારે
જ્યારે $\Gamma \Entails !A$ હોય ત્યારે $!A \in \Gamma$. !!{sentence}sના
ગણ~$\Gamma$નું \emph{સંવરણ} એટલે
$\Setabs{!A}{\Gamma \Entails !A}$.

જો~$\Gamma$, વાક્યોના ગણ~$\Delta$નું સંવરણ હોય, તો આપણે કહીએ કે
$\Gamma$ને~$\Delta$એ \emph{સ્વયંસિદ્ધીકૃત} કર્યો છે.
\end{defn}

\begin{explain}
સ્વયંસિદ્ધિમૂલક સિદ્ધાંતને તેના સ્વયંસિદ્ધિઓના ગણ~$\Delta$ દ્વારા
સ્વયંસિદ્ધીકૃત વાક્યોના ગણ તરીકે વિચારી શકીએ છીએ. બીજા શબ્દોમાં,
આપણે જે સ્વયંસિદ્ધિમૂલક રીતે વિકસિત વિજ્ઞાનનો અભ્યાસ કરવા માગીએ તેના
પાયાના ખ્યાલો માટેના અતાર્કિક સંકેતો ધરાવતી પ્રથમ-ક્રમ ભાષા અને એ
વિજ્ઞાનના પાયાના નિયમો વ્યક્ત કરતાં !!{sentence}sનો ગણ આપેલો હોય, તો
સ્વયંસિદ્ધિઓમાંથી ફલિત થતાં આ ભાષાનાં બધાં !!{sentence}s દ્વારા
પ્રતિનિધિત સિદ્ધાંતને વિચારી શકીએ છીએ. તેમાં માત્ર એક પાયાનો ખ્યાલ અને
સરળ સ્વયંસિદ્ધિઓ ધરાવતા આંશિક ક્રમોના સિદ્ધાંત જેવા સરળ દાખલાથી લઈને
ન્યૂટનીય યંત્રવિજ્ઞાન જેવા જટિલ સિદ્ધાંતો સુધીનો સમાવેશ થાય છે.

સ્વયંસિદ્ધિમૂલક પદ્ધતિ પ્રત્યેના આ વિધિવત્ અભિગમને એટલો મહત્ત્વનો
બનાવતી મુખ્ય તાર્કિક હકીકતો નીચેની છે. ધારો કે~$\Gamma$ કોઈ સિદ્ધાંત
માટેનું સ્વયંસિદ્ધિ-તંત્ર, એટલે કે વાક્યોનો ગણ, છે.
\begin{enumerate}
\item સ્વયંસિદ્ધિ-તંત્ર અભિપ્રેત !!{structure}sના વર્ગને ક્યારે ચોક્કસ
  રીતે વર્ણવે છે તે આપણે કહી શકીએ છીએ. એટલે કે, જો આપણને
  !!{structure}sના કોઈ ચોક્કસ વર્ગમાં રસ હોય, તો સ્વયંસિદ્ધિ-તંત્ર
  ~$\Gamma$ એ વર્ગને ત્યારે અને માત્ર ત્યારે જ સફળતાથી વર્ણવે છે જ્યારે
  એ !!{structure}s બરાબર એવી~$\Struct M$ હોય કે $\Sat{M}{\Gamma}$.
\item આ બાબતમાં આપણે નિષ્ફળ જઈ શકીએ, કારણ કે એવી~$\Struct M$ હોઈ શકે
  કે $\Sat{M}{\Gamma}$, પરંતુ~$\Struct M$ આપણી અભિપ્રેત
  !!{structure}sમાંની ન હોય. તેથી કદાચ આપણે~$\Struct M$માં અસત્ય હોય
  એવી સ્વયંસિદ્ધિઓ ઉમેરવી પડે.
\item ઓછામાં ઓછું~$\Gamma$ બધી અભિપ્રેત !!{structure}sમાં સત્ય છે એ
  બાબતમાં આપણે સફળ હોઈએ, તો જ્યારે પણ $\Gamma \Entails !A$, ત્યારે
  વાક્ય~$!A$ બધી અભિપ્રેત !!{structure}sમાં સત્ય હોય છે. તેથી વાક્યો
  સ્વયંસિદ્ધિઓમાંથી ફલિત થાય છે એટલું બતાવીને જ તેઓ બધી અભિપ્રેત
  !!{structure}sમાં સત્ય છે એવું દર્શાવવા આપણે તાર્કિક સાધનો (જેમ કે
  !!{derivation} પદ્ધતિઓ) વાપરી શકીએ છીએ.
\item કેટલીક વાર આપણા મનમાં અભિપ્રેત !!{structure}s હોતી નથી; તેના
  બદલે આપણે સ્વયંસિદ્ધિઓથી જ શરૂ કરીએ છીએ: કેટલાક પાયાના ખ્યાલો લઈએ
  છીએ અને તેઓ અમુક નિયમોને સંતોષે એમ માગીએ છીએ; પછી એ નિયમોને
  સ્વયંસિદ્ધિ-તંત્રમાં સંહિતાબદ્ધ કરીએ છીએ. આપણે તરત ચકાસવા માગીએ એવી
  એક બાબત એ છે કે સ્વયંસિદ્ધિઓ પરસ્પર વિસંવાદી નથી: જો હોય, તો આ
  નિયમોનું પાલન કરતા કોઈ ખ્યાલ હોઈ શકે નહિ અને આપણે અસંગત સિદ્ધાંત
  રચવાનો પ્રયત્ન કર્યો છે. $\Gamma$નું એક નિદર્શ શોધીને આવું થતું નથી
  તે ચકાસી શકીએ છીએ. આપણા સિદ્ધાંતનાં નિદર્શો હોય, તો તાર્કિક પદ્ધતિઓ
  વડે તેમનો અભ્યાસ અને નિદર્શોની રચના બંને કરી શકીએ છીએ.
\item સ્વયંસિદ્ધિઓની સ્વતંત્રતા પણ એટલો જ અગત્યનો પ્રશ્ન છે. કોઈ એક
  સ્વયંસિદ્ધિ બીજાંનું ફલિત હોઈ શકે છે અને તેથી અનાવશ્યક હોઈ શકે છે.
  $\Gamma$માંનું સ્વયંસિદ્ધિ~$!A$ અનાવશ્યક છે એ
  $\Gamma \setminus \{!A\} \Entails !A$ સાબિત કરીને બતાવી શકીએ છીએ.
  વળી, $(\Gamma \setminus \{!A\}) \cup \{\lnot !A\}$ સંતોષ્ય છે એમ
  બતાવીને તે સ્વયંસિદ્ધિ અનાવશ્યક નથી એ પણ સાબિત કરી શકીએ છીએ. દાખલા
  તરીકે, ભૂમિતિની બીજી સ્વયંસિદ્ધિઓથી સમાંતર સ્વયંસિદ્ધિ સ્વતંત્ર છે
  તે આ જ રીતે બતાવવામાં આવ્યું હતું.
\item બીજો અગત્યનો પ્રશ્ન સિદ્ધાંતમાંના ખ્યાલોની વ્યાખ્યેયતાનો છે:
  ભાષાની પસંદગી સિદ્ધાંતનાં નિદર્શો શેનાં બને છે તે નક્કી કરે છે.
  પરંતુ સિદ્ધાંતના દરેક પાસાને તેના નિદર્શોમાં અલગથી રજૂ કરવો જરૂરી
  નથી. દાખલા તરીકે, દરેક ક્રમ~$\le$ તેને અનુરૂપ કડક ક્રમ~$<$ નક્કી કરે
  છે---એક આપેલો હોય તો બીજાને વ્યાખ્યાયિત કરી શકીએ છીએ. તેથી આવા ક્રમ
  ધરાવતા સિદ્ધાંતના નિદર્શમાં તેને અનુરૂપ કડક ક્રમ પણ અલગથી હોવો
  જરૂરી નથી. સામાન્ય રીતે, કોઈ એક સંબંધને બીજા સંબંધોના પદોમાં ક્યારે
  વ્યાખ્યાયિત કરી શકાય? કોઈ સંબંધને બીજા સંબંધોના પદોમાં વ્યાખ્યાયિત
  કરવો ક્યારે અશક્ય હોય (અને તેથી તેને ભાષાના પાયાના સંકેતોમાં ઉમેરવો
  પડે)?
\end{enumerate}
\end{explain}


\end{document}
